% _highlight.tex
% Macros for visually marking .tex updates between revisions.
% Usage during an edit:
%   \diffold{previous statement of Theorem 2}
%   \diffnew{revised statement of Theorem 2}
% Before the NEXT edit pass: remove \diffold{...} entirely, and unwrap
% \diffnew{...} down to its plain content (drop the macro, keep the text).

\definecolor{difnewcolor}{RGB}{0,102,204}   % new/changed content
\definecolor{difoldcolor}{RGB}{150,150,150} % superseded content

% Declared with +m so a highlighted block may span several paragraphs;
% \textcolor is short, and a \par inside its argument aborts the run.
% A run-in heading such as \paragraph is typeset only when its paragraph
% starts. \diffstartpar starts that paragraph before the color is set, so the
% heading keeps the text color; elsewhere (\if@noskipsec false) it does nothing.
\makeatletter
\newcommand{\diffstartpar}{\if@noskipsec\leavevmode\fi}
\makeatother
\NewDocumentCommand{\diffnew}{+m}{\diffstartpar{\color{difnewcolor}#1}}
% \diffold strikes its argument through with \sout (ulem), which can only set
% a run of words. The argument is therefore handled in pieces:
%  * it is split at paragraph breaks, and each paragraph is struck separately,
%    so the paragraphs of a superseded block stay paragraphs;
%  * a paragraph that starts with a heading (\section ... \subparagraph) is
%    issued as a real heading with its title struck through, and the rest of
%    the paragraph is struck as usual; the table of contents, running header,
%    and bookmarks get the plain title;
%  * \sout cannot strike \cite or \parencite once biblatex has resolved them
%    ("Extra }, or forgotten \endgroup"), but it can strike a box, so every
%    citation command, with its optional arguments and keys, is put in \mbox.
% Lists, display math, and headings in the middle of a paragraph are still
% beyond \sout.
\NewDocumentCommand{\diffoldtitle}{m}{\sout{#1}}
\ExplSyntaxOn
\tl_new:N \l__diff_old_tl
\seq_new:N \l__diff_old_seq
\bool_new:N \l__diff_first_bool
\bool_new:N \l__diff_heading_bool
\regex_const:Nn \c__diff_cite_regex
  {
    \c{(?:cite|Cite|parencite|Parencite|textcite|Textcite|autocite|Autocite
         |citeauthor|Citeauthor|citeyear)}
    \*? (?: \x{5B} [^\x{5D}]* \x{5D} ){0,2} \cB. (?: \c[^BE]. )* \cE.
  }
\NewDocumentCommand{\diffold}{+m}
  {
    \diffstartpar
    \group_begin:
      \color{difoldcolor}
      \tl_set:Nn \l__diff_old_tl {#1}
      \regex_replace_all:NnN \c__diff_cite_regex
        { \c{mbox} \cB\{ \0 \cE\} } \l__diff_old_tl
      \seq_set_split:NnV \l__diff_old_seq { \par } \l__diff_old_tl
      \bool_set_true:N \l__diff_first_bool
      \seq_map_inline:Nn \l__diff_old_seq
        {
          \tl_if_blank:nF {##1}
            {
              \bool_if:NF \l__diff_first_bool { \par }
              \bool_set_false:N \l__diff_first_bool
              \__diff_old_paragraph:n {##1}
            }
        }
    \group_end:
  }
% One paragraph: a leading heading command is issued by \__diff_old_heading:,
% everything else is struck through.
\cs_new_protected:Npn \__diff_old_paragraph:n #1
  {
    \bool_set_false:N \l__diff_heading_bool
    \clist_map_inline:nn
      { \section, \subsection, \subsubsection, \paragraph, \subparagraph }
      {
        \tl_if_head_eq_meaning:nNT {#1} ##1
          { \bool_set_true:N \l__diff_heading_bool \clist_map_break: }
      }
    \bool_if:NTF \l__diff_heading_bool
      { \__diff_old_heading: #1 \q__diff_stop }
      { \sout {#1} }
  }
\quark_new:N \q__diff_stop
% <heading command> <*> [<short title>] {<title>} <rest of the paragraph>
\NewDocumentCommand{\__diff_old_heading:}{m s o m u{\q__diff_stop}}
  {
    \IfBooleanTF {#2}
      { #1 * { \diffoldtitle {#4} } }
      { #1 [ \IfValueTF {#3} {#3} {#4} ] { \diffoldtitle {#4} } }
    \diffstartpar
    \tl_set:Nn \l__diff_old_tl {#5}
    \tl_trim_spaces:N \l__diff_old_tl
    \tl_if_empty:NF \l__diff_old_tl { \exp_args:NV \sout \l__diff_old_tl }
  }
\ExplSyntaxOff

% Display-math variants: \sout cannot strike a display equation directly, but it
% strikes a box, so the equation is boxed first.  Use \diffoldmath inside an
% equation*/align* body for a superseded formula, and \diffoldmathwide when the
% formula would otherwise overflow the text width.
\newcommand{\diffoldmath}[1]{%
  \textcolor{difoldcolor}{\sout{\mbox{\ensuremath{\displaystyle #1}}}}%
}
\newcommand{\diffoldmathwide}[1]{%
  \textcolor{difoldcolor}{%
    \resizebox{\textwidth}{!}{\sout{\mbox{\ensuremath{\displaystyle #1}}}}%
  }%
}
